\section{الحركيات الأمامية}

تُستخدم الكينماتيكا المباشرة \LR{(Forward Kinematics)} لتحديد موضع واتجاه نهاية الذراع الروبوتية \LR{(End-Effector)} اعتماداً على قيم زوايا المفاصل. تُعد هذه العملية أساسية في تحليل حركة الروبوت، حيث تسمح بربط فضاء المفاصل \LR{(Joint Space)} بفضاء العمل \LR{(Workspace)}.

في هذا المشروع، تم استخدام معاملات دينافيت–هارتنبرغ \LR{(Denavit--Hartenberg)} لاشتقاق مصفوفات التحويل المتجانس بين الإطارات المتتالية. حيث تُعطى مصفوفة التحويل بين الإطارين \(i-1\) و\(i\) بالشكل العام:

\[
A_i =
\begin{bmatrix}
\cos\theta_i & -\sin\theta_i \cos\alpha_i & \sin\theta_i \sin\alpha_i & a_i \cos\theta_i \\
\sin\theta_i & \cos\theta_i \cos\alpha_i  & -\cos\theta_i \sin\alpha_i & a_i \sin\theta_i \\
0            & \sin\alpha_i              & \cos\alpha_i               & d_i \\
0            & 0                         & 0                          & 1
\end{bmatrix}
\]

تمثل هذه المصفوفة الدوران والإزاحة بين الإطارين المتتاليين، حيث تعتمد على معاملات DH لكل وصلة.

بالتطبيق على الذراع الروبوتية المدروسة، وبالاعتماد على جدول معاملات دينافيت–هارتنبرغ الذي تم استخراجه سابقاً، يمكن الحصول على مصفوفات التحويل المتجانس الخاصة بكل وصلة كما يلي:

\subsection{مصفوفات التحويل}


\[
T_1^0 =
\begin{bmatrix}
\cos\theta_{1} & 0 & \sin\theta_{1}  & 0.017\cos\theta_{1} \\
\sin\theta_{1} & 0 & -\cos\theta_{1} & 0.017\sin\theta_{1} \\
0              & 1 &               0 & 0.05 \\
0              & 0 &               0 & 1
\end{bmatrix}
\quad 
T_2^1 =
\begin{bmatrix}
-\sin\theta_2 & -\cos\theta_2 & 0 & -0.12\sin\theta_2 \\
\cos\theta_2 & -\sin\theta_2 & 0 & 0.12\cos\theta_2 \\
0 & 0 & 1 & 0 \\
0 & 0 & 0 & 1
\end{bmatrix}
\]



\[
T_3^2 =
\begin{bmatrix}
\cos\theta_{3} & 0 & \sin\theta_{3} & 0.044\cos\theta_{3} \\
\sin\theta_{3} & 0 & -\cos\theta_{3} & 0.044\sin\theta_{3} \\
0 & 1 & 0 & 0 \\
0 & 0 & 0 & 1
\end{bmatrix}
\quad
T_4^3 =
\begin{bmatrix}
-\cos\theta_{4} & 0 & -\sin\theta_{4} & 0 \\
-\sin\theta_{4} & 0 & \cos\theta_{4}  & 0 \\
0               & 1 & 0               & 0.15 \\
0               & 0 & 0               & 1
\end{bmatrix}
\]

\[
\makebox[\textwidth][c]{
$
T_5^4 =
\begin{bmatrix}
-\sin\theta_{5} & -\cos\theta_{5} & 0 & -0.152\sin\theta_{5} \\
\cos\theta_{5}  & -\sin\theta_{5} & 0 & 0.152\cos\theta_{5} \\
0               & 0               & 1 & 0 \\
0               & 0               & 0 & 1
\end{bmatrix}
$
}
\]






حيث تمثل كل مصفوفة من هذه المصفوفات العلاقة بين إطارين متتاليين في الذراع الروبوتية.

وبالتالي، يتم الحصول على مصفوفة التحويل الكلية من القاعدة إلى نهاية الذراع من خلال:

\[
T_0^5 = T_0^1 \cdot T_{12} \cdot T_2^3 \cdot T_3^4 \cdot T_4^5
\]

والتي تعبّر عن موضع واتجاه نهاية الذراع الروبوتية بالنسبة لإطار القاعدة.

من خلال هذه المصفوفة، يمكن استخراج:
\begin{itemize}
    \item موضع نهاية الذراع (Position)
    \item اتجاه نهاية الذراع (Orientation)
\end{itemize}

وبالتالي، تُمكّن الكينماتيكا المباشرة من توصيف الحركة الكاملة للذراع الروبوتية انطلاقاً من زوايا المفاصل المعطاة.